% SOURCE_AUTHORITY: sga5_fr_workpass.tex, lines 12434--12472 (LF-normalized unit).
% SOURCE_SHA256: 791F4EFFC5E02832D5D77ED03518C8156D6F07E4C8238B03545DB93D883FBB28
\subsubsection*{3.7}

Sea $P$ un $\Lambda[G]$-módulo. Es, en particular, un $\Lambda$-módulo. Puede dotarse al $\Lambda$-módulo $\check P=\Hom_\Lambda(P,\Lambda)$ de una estructura de $\Lambda[G]$-módulo por la izquierda poniendo, para cada $u\in\Hom_\Lambda(P,\Lambda)$ y $g\in G$, $(gu)(x)=u(g^{-1}x)$, es decir,
\[
g_{\check P}={}^{t}(g^{-1})_P.
\tag{3.7.1}
\]
El módulo $\check P$ sobre $\Lambda[G]$ así obtenido se llamará el dual de $P$.

Puede definirse, por tanto, el dual $\check P^\bullet=\Hom_\Lambda^\bullet(P^\bullet,\Lambda)$ de un complejo $P^\bullet$ de $\Lambda[G]$-módulos como un complejo de $\Lambda[G]$-módulos.

Si $\chi_P$ es el carácter de $P$, el carácter $\chi_{\check P}$ de su dual viene dado por la fórmula:
\[
\chi_{\check P}(g)=\chi_P(g^{-1}),
\tag{3.7.2}
\]
que se deduce inmediatamente de (4.3.1) utilizando la fórmula:
\[
\Tr_\Lambda^*(v)=\Tr_\Lambda^*({}^{t}v).
\]
Si $P$ es un $\Lambda[G]$-módulo proyectivo de tipo finito, es claro que $\check P$ es también un $\Lambda[G]$-módulo proyectivo de tipo finito. Además, en este caso, el homomorfismo canónico de $\Lambda$-módulos
\[
P\longrightarrow\check{\check P}
\]
es, de hecho, un isomorfismo de $\Lambda[G]$-módulos. De este modo se define sobre $K^{\bullet}(\Lambda[G])$ una involución
\[
x\longmapsto \check x,
\]
que asocia a la clase de $P$ la clase de $\check P$. Es inmediato que, si $\mu:\Lambda\to\Gamma$ es un homomorfismo de anillos conmutativos y $u:\Lambda[G]\to\Gamma[G]$ el homomorfismo asociado, el homomorfismo
\[
u^*:K^{\bullet}(\Lambda[G])\longrightarrow K^{\bullet}(\Gamma[G])
\]
conmuta con las involuciones $x\longmapsto \check x$ sobre $K^{\bullet}(\Lambda[G])$ y $K^{\bullet}(\Gamma[G])$.

Sea $A$ un álgebra sobre $\Lambda$, y sea
\[
A[G]=\Lambda[G]\otimes_\Lambda A.
\]
Dado un complejo $X^\bullet$ de $A[G]$-módulos, se denota por $X^{\bullet G}$ el complejo de $A$-módulos imagen de $X^\bullet$ por el funtor que asocia a cada $A[G]$-módulo $X$ el $A$-módulo formado por los elementos de $X$ invariantes bajo $G$.
